\zhbox{54}{67}{504}{128}{10}{\textbf{USB Power Delivery（PD）应用的 I\textsuperscript{2}C VREF 设置顺序}\par
拔出受电设备时，PD 控制器将 ENPWR 设为0以禁用总线电压（$V_{BUS}$）。重新接入受电设备时，PD 控制器按以下顺序将 $V_{BUS}$ 设为5 V：\par
1. 根据新的 $V_{OUT}$ 写入 VREF 位。\par
2. 写入 GO\_BIT=1，并保持 ENPWR=0，使 $V_{REF}$ 达到目标值。\par
3. 写入 ENPWR=1，使能输出。\par
图11显示新的 $V_{OUT}$ 请求顺序。}
\zhbox{218}{232}{61}{24}{8}{请求新的 $V_{OUT}$}
\zhbox{299}{275}{106}{14}{8}{2. 写入 GO\_BIT=1}
\zhbox{247}{325}{114}{32}{8}{1. 根据新的请求写入 VREF 寄存器}
\zhbox{366}{324}{140}{32}{8}{3. 写入 ENPWR=1}
\zhbox{54}{361}{504}{16}{10}{\centering\textbf{图11：新输出电压请求顺序}}
\zhbox{54}{379}{504}{15}{10}{\textbf{寄存器0x03：CONTROL\_2（读/写）}}
\zhbox{55}{397}{40}{11}{8}{\centering 位}
\zhbox{99}{397}{54}{11}{8}{\centering 名称}
\zhbox{157}{397}{44}{11}{8}{\centering 默认值}
\zhbox{207}{397}{350}{11}{8}{说明}
\zhbox{207}{415}{348}{79}{8.5}{设置开关频率（$f_{SW}$）变化率。FSW 默认值由启动时 FS 引脚设置决定；各电平设置见第9页电气特性。\\
00：200 kHz变化率\\01：300 kHz变化率\\10：400 kHz变化率\\11：600 kHz变化率}
\zhbox{207}{501}{348}{36}{8.5}{设置升降压开关频率。\\0：设为基本开关频率的40\%。\\1：设为基本开关频率的80\%。}
\zhbox{207}{545}{348}{90}{8.5}{设置降压模式电流达到逐周期谷值限值（$I_{LIMIT\_VALLEY}$）或升压模式达到逐周期峰值限值（$I_{LIMIT\_PEAK}$）时的模式。\par
00：无打嗝或锁断保护；逐周期限制电感电流（$I_L$）。\\
01：电流达到 $I_{LIMIT}$ 且 $V_{FB}<60\%V_{REF}$ 时进入打嗝模式；$t_{SS}$ 由 SS 放电通路控制。\\
10：触发 OCP；器件通过软启动恢复正常工作。\\11：保留。}

